% BEGIN R28_BOUNDARY_OCCUPANCY
% Insertar en MG01 tras la definición de c_j,o_j, antes de su valor.
La ley de la frontera de profundidad seis fija \(c_j=1\). Con el eje \(z\) y el
agregado \(q\) ya determinados, las sumas enteras laterales son
\[
 (x_j+y_j)_{j=1}^6=(2,4,3,4,4,2).
\]
La regla de ocupación mínima selecciona, entre sus descomposiciones con
\(x_j,y_j\in\{0,1,2\}\), las columnas con menor número de entradas
no nulas, contando también \(z_j\). Para suma dos, las parejas extremas
\((0,2),(2,0)\) ocupan una posición lateral; para suma tres, las parejas
\((1,2),(2,1)\) ocupan dos; para suma cuatro, sólo existe \((2,2)\).
Al añadir la ocupación del eje se obtiene \(o=(2,2,3,2,3,2)\).
En esta frontera, con \(s=[x+y]_3\), la misma regla se expresa como
\[
 o_j=2+\mathbf1_{\{z_j\ne0,\ s_j\ne2\}}.
\]
% END R28_BOUNDARY_OCCUPANCY
